% !TeX root = ../../Book1.tex
% 纲要标题：### 4.2 共轭作用
% 纲要位置：计划纲要.md，第 223 行。

\section{共轭作用}
\label{b1:sec:0402}

上一节允许群作用在任意集合上，现在取这个集合为群自身。左乘把所有元素连成一个轨道，共轭却保留了更细的区别：单位始终不动，其他元素则按共轭类分组。计算这些轨道的大小，可以把群阶分解成若干子群指数之和；对素数幂阶群，这一等式尤其有力。

% 本节正文按下列原纲要要点撰写。

% 纲要内容（仅作写作计划，尚非教材正文）：
%
% * 共轭类、中心化子与正规化子
% * 类方程
% * 有限 p 群的中心
%

\subsection{共轭类、中心化子与正规化子}
\label{b1:subsec:0402:01}

令 \(g\cdot x=gxg^{-1}\)，则 \(1\cdot x=x\)，并且
\[
g\cdot(h\cdot x)=ghx(gh)^{-1}=(gh)\cdot x,
\]
故这确为群作用。

\begin{definition}\label{b1:def:centralizer-normalizer}
设 \(G\) 为群，\(x\in G\)。将 \(x\) 在 \(G\) 的共轭作用下的轨道称为它的\term[gonge-lei]{共轭类}[conjugacy class]，记为
\(\operatorname{Cl}_G(x)=\{\,gxg^{-1}\mid g\in G\,\}\)。
元素 \(x\) 在该作用下的稳定子群称为它的\term[zhongxinhua-zi]{中心化子}[centralizer]，记为
\[
C_G(x)\coloneqq\{\,g\in G\mid gx=xg\,\}.
\]
对任意子集 \(A\subseteq G\)，其中心化子为
\[
C_G(A)\coloneqq\{\,g\in G\mid ga=ag\text{ 对每个 }a\in A\,\}
=\bigcap_{a\in A}C_G(a).
\]
前面定义的中心可以写为 \(Z(G)=C_G(G)\)。
子群 \(H\leq G\) 的\term[zhengguihua-zi]{正规化子}[normalizer]为
\[
N_G(H)\coloneqq\{\,g\in G\mid gHg^{-1}=H\,\}.
\]
\symidx{\(\operatorname{Cl}_G(x)\)：元素的共轭类}
\symidx{\(C_G(x)\)：元素的中心化子}
\symidx{\(C_G(A)\)：子集的中心化子}
\symidx{\(N_G(H)\)：子群的正规化子}
\end{definition}

这里 \(C_G(x)\) 固定的是元素 \(x\)。若考察整个共轭类作为集合是否保持不变，答案则不同：对任意 \(g\in G\)，都有
\[
g\operatorname{Cl}_G(x)g^{-1}
=\Bigl\{\,(gh)x(gh)^{-1}\mid h\in G\,\Bigr\}
=\operatorname{Cl}_G(x),
\]
因为 \(h\mapsto gh\) 是 \(G\) 上的双射。因此，保持这个共轭类整体不变的群元素是 \(G\) 的全部元素，而不只是 \(C_G(x)\) 中的元素。

\ProseParagraphBreak
正规化子对应于另一个作用，其作用集合是 \(G\) 的全部子群组成的集合 \(\mathcal S(G)\)。对 \(g\in G\) 和 \(H\in\mathcal S(G)\)，定义
\[
g\cdot H=gHg^{-1}.
\]
这首先确实给出一个子群：\(1\in gHg^{-1}\)，而对 \(u,v\in H\)，有
\[
(gug^{-1})(gvg^{-1})^{-1}=g(uv^{-1})g^{-1}\in gHg^{-1},
\]
故\cref{b1:prop:subgroup-test}适用。其次，\(1\cdot H=H\)，并且
\[
g\cdot(h\cdot H)=g(hHh^{-1})g^{-1}=(gh)H(gh)^{-1}=(gh)\cdot H.
\]
因此这定义了 \(G\) 在 \(\mathcal S(G)\) 上的作用。点 \(H\) 的稳定子群恰为 \(N_G(H)\)，所以正规化子是子群。由\cref{b1:thm:orbit-stabilizer}的陪集与轨道双射，还得到
\[
\bigl|\{\,gHg^{-1}\mid g\in G\,\}\bigr|=[G:N_G(H)].
\]
此处不要求 \(G\) 有限；等号表示共轭子群集合与相应左陪集集合具有相同基数。

\ProseParagraphBreak
中心化子 \(C_G(H)\) 要求共轭固定 \(H\) 的每个元素，正规化子 \(N_G(H)\) 则只要求共轭保持 \(H\) 整体不变。每个 \(h\in H\) 都满足 \(hHh^{-1}=H\)，故 \(H\leq N_G(H)\)；再按正规化子的定义，便有 \(H\trianglelefteq N_G(H)\)。此外，\(H\trianglelefteq G\) 当且仅当 \(N_G(H)=G\)。中心也是正规子群：\(Z(G)=\bigcap_{x\in G}C_G(x)\) 是子群，而且若 \(z\in Z(G)\)，则 \(gzg^{-1}=z\)，故每个共轭仍在中心中。

\ProseParagraphBreak
以 \(S_3\) 为例，\(C_{S_3}\Bigl((12)\Bigr)=\Bigl\{1,(12)\Bigr\}\)，因为与 \((12)\) 交换的置换必须保持其唯一固定点 \(3\)，并只能在 \(1,2\) 之间置换。因此换位的共轭类有三个元素。三个换位组成的集合被整个 \(S_3\) 保持，而其中单个换位的中心化子只有两个元素。另一方面，\(H=\Bigl\langle(123)\Bigr\rangle\) 的中心化子为 \(H\)，而正规化子是整个 \(S_3\)；换位把 \((123)\) 变成其逆，保持子群，却不逐点固定子群。

\subsection{类方程}
\label{b1:subsec:0402:02}

对共轭作用使用\cref{b1:thm:orbit-stabilizer}，其中点 \(x\) 的稳定子群为 \(C_G(x)\)，得到
\[
\bigl|\operatorname{Cl}_G(x)\bigr|=\bigl[G:C_G(x)\bigr].
\]
此式对任意群都成立，因为该定理给出了两个集合之间的双射。共轭类大小为 \(1\) 当且仅当 \(x\in Z(G)\)。以下设 \(G\) 有限，把这些单元素轨道先分离出来，再把其他共轭类的大小相加，就得到\term[lei-fangcheng]{类方程}[class equation]。

\begin{proposition}[类方程]\label{b1:prop:class-equation}
设 \(G\) 为有限群，从每个非中心共轭类中选一个代表 \(x_1,\ldots,x_r\)。则
\begin{equation}\label{b1:eq:class-equation}
|G|=\bigl|Z(G)\bigr|+\sum_{i=1}^{r}\bigl[G:C_G(x_i)\bigr].
\end{equation}
\end{proposition}
\begin{proof}
共轭类把 \(G\) 分割成不交并。中心元素各自构成一个单元素类。对每个非中心代表 \(x_i\)，\cref{b1:thm:orbit-stabilizer}给出其共轭类大小为 \(\bigl[G:C_G(x_i)\bigr]\)。对这些不交的类计数，就得到\eqref{b1:eq:class-equation}。
\end{proof}

例如 \(S_3\) 的类方程为 \(6=1+3+2\)：单位、三个换位、两个三轮换恰好给出三个共轭类。此式中的 \(1\) 表明 \(Z(S_3)=\{1\}\)。类方程中的求和按共轭类取代表，而不是对每个非中心元素求和；否则同一个轨道会重复计算。

\ProseParagraphBreak
对 \(D_8=\langle r,s\rangle\)，其中 \(r\) 是正方形四分之一周旋转，\(s\) 是反射，有 \(srs^{-1}=r^{-1}\)。旋转 \(r^a\) 与 \(s\) 交换当且仅当 \(r^a=r^{-a}\)，即 \(a\) 为偶数；反射 \(r^as\) 不与 \(r\) 交换。因此 \(Z(D_8)=\{1,r^2\}\)。由同一关系可得
\[
sr^as^{-1}=r^{-a},\qquad
r^b(r^as)r^{-b}=r^{a+2b}s,\qquad
s(r^as)s^{-1}=r^{-a}s.
\]
旋转彼此交换，而用 \(s\) 共轭将旋转变成其逆，所以非中心旋转恰组成 \(\{r,r^3\}\)。对反射，用生成元 \(r,s\) 共轭都保持指数的奇偶性，故任意群元素的共轭也保持这一奇偶性；用 \(r\) 共轭又将指数加 \(2\)，从而遍历同一奇偶类。于是全部共轭类为
\[
\{1\},\quad \{r^2\},\quad \{r,r^3\},\quad
\{s,r^2s\},\quad \{rs,r^3s\},
\]
类方程为 \(8=2+2+2+2\)。它既反映中心，也区分了正方形两种几何位置不同的反射。

\subsection{有限 \texorpdfstring{\(p\)}{p} 群的中心}
\label{b1:subsec:0402:03}

\begin{definition}\label{b1:def:p-group}
设 \(p\) 为素数，\(a\geq0\) 为整数。如果有限群 \(G\) 的阶为 \(p^a\)，则将 \(G\) 称为有限 \(p\)-群。当 \(a\geq1\) 时，\(G\) 非平凡。
\end{definition}

由\cref{b1:thm:orbit-stabilizer}，有限 \term[p-qun]{\(p\)-群}[p-group] \(P\) 的每个轨道大小为某个稳定子群的指数 \([P:P_x]\)，因而整除 \(|P|\)，是 \(p\) 的幂。这给出一个将在 Sylow 理论中反复使用的计数方法。

\ProseParagraphBreak
对于任意群 \(G\) 在集合 \(X\) 上的作用，把被所有群元素固定的点称为\term[gongtong-budongdian]{共同不动点}[common fixed point]，这些点组成的集合记为
\[
X^G\coloneqq\{\,x\in X\mid gx=x\text{ 对每个 }g\in G\,\}.
\]
\symidx{\(X^G\)：群作用的共同不动点集}
共同不动点恰是单元素轨道中的点。作用群为 \(P\) 时，这个集合相应地记为 \(X^P\)。
\symidx{\(X^P\)：\(p\)-群作用的共同不动点集}

\begin{lemma}\label{b1:lem:p-action-fixed}
设 \(p\) 为素数。若有限 \(p\)-群 \(P\) 作用在有限集合 \(X\) 上，记
\(X^P=\{\,x\in X\mid gx=x\text{ 对所有 }g\in P\,\}\)，则
\[
|X|\equiv |X^P|\pmod p.
\]
\end{lemma}
\begin{proof}
大小为 \(1\) 的轨道恰由 \(X^P\) 中的点组成；其他轨道大小是大于 \(1\) 的 \(p\) 幂，均被 \(p\) 整除。把轨道大小相加再模 \(p\) 即得结论。
\end{proof}
\begin{theorem}\label{b1:thm:p-group-center}
设 \(p\) 为素数，\(P\) 为非平凡有限 \(p\)-群。则 \(Z(P)\) 非平凡，而且 \(|Z(P)|\) 被 \(p\) 整除。
\end{theorem}
\begin{proof}
记所给非平凡有限 \(p\)-群为 \(P\)，令它作用在有限集合 \(X=P\) 上，作用为 \(g\cdot x=gxg^{-1}\)。一个元素 \(x\) 被所有 \(g\in P\) 固定，当且仅当它与所有 \(g\) 交换，所以 \(X^P=Z(P)\)。对这个作用应用\cref{b1:lem:p-action-fixed}，其结论 \(|X^P|\equiv|X|\pmod p\) 在这里成为
\[
\bigl|Z(P)\bigr|\equiv|P|\equiv0\pmod p.
\]
末一个同余使用 \(P\neq\{1\}\)，即 \(|P|=p^a\)、\(a\geq1\)。中心包含单位，故其阶是正的 \(p\) 倍数，至少为 \(p\)，从而中心非平凡。
\end{proof}

这里“不动点存在”并非仅由作用非空得到，而是由同余和中心含单位共同得到。一般有限群的中心可以平凡，如 \(S_3\) 所示。

\ProseParagraphBreak
下面先记录一个不要求有限性的事实，供后续讨论中心商时使用。

\begin{proposition}\label{b1:prop:cyclic-central-quotient}
若群 \(G\) 的商群 \(G/Z(G)\) 是循环群，则 \(G\) 是交换群。
\end{proposition}
\begin{proof}
取一个生成陪集 \(aZ(G)\)。对每个 \(g\in G\)，存在整数 \(i\) 使 \(gZ(G)=a^iZ(G)\)，故 \(g=a^iz\)，其中 \(z\in Z(G)\)。任取两个元素 \(a^iz,a^jw\)，利用 \(z,w\) 与所有元素交换，得到
\[
(a^iz)(a^jw)=a^{i+j}zw=a^{j+i}wz=(a^jw)(a^iz).
\]
因此 \(G\) 中任意两个元素交换。
\end{proof}

\begin{corollary}\label{b1:cor:order-p-square-abelian}
设 \(p\) 为素数。每个阶为 \(p^2\) 的群都是交换群。
\end{corollary}
\begin{proof}
由\cref{b1:thm:p-group-center}，\(Z(G)\) 非平凡；又由\cref{b1:thm:lagrange}，其阶整除 \(p^2\)，所以只能为 \(p\) 或 \(p^2\)。若为 \(p^2\)，则 \(Z(G)=G\)，结论成立。若为 \(p\)，商群 \(G/Z(G)\) 的阶为 \(p\)，由\cref{b1:cor:finite-order-divides}中关于素数阶群的结论，它是循环群。再用\cref{b1:prop:cyclic-central-quotient}便知 \(G\) 交换；事实上这也排除了中心阶恰为 \(p\) 的可能。
\end{proof}

中心非平凡并不意味着整个 \(p\)-群交换；阶为 \(8\) 的 \(D_8\) 和 \(Q_8\) 已提供反例。上述推论利用的是商群阶恰为素数，不能把这一论证无条件移到 \(p^3\) 阶。
